\chapter{QML desde cero}
\label{cap:qml}

QML es el lenguaje con el que Qt describe interfaces. Si vienes de HTML/CSS/JavaScript te
resultará familiar: es \emph{declarativo} (describes qué hay en pantalla) y la lógica se
escribe en JavaScript. Si vienes de Java Swing o de Android, la diferencia principal es que en
QML casi nunca «actualizas» la pantalla a mano: declaras de qué depende cada cosa y Qt la
recalcula sola.

\section{Un archivo QML}

Cada archivo \archivo{qml/Nombre.qml} define un \emph{componente} que se usa en otros archivos
como \cod{Nombre { ... }}. Dentro hay un árbol de objetos, cada uno con propiedades:

\begin{qml}
import QtQuick
import app.musicplayer        // el módulo del proyecto: Theme, Player, Cover...

Rectangle {
    id: boton                 // nombre para referirse a este objeto desde otros
    width: 160
    height: 44
    radius: height / 2        // binding: depende de height
    color: Theme.control

    Text {
        anchors.centerIn: parent         // colócate en el centro de tu padre
        text: "Reproducir"
        color: Theme.text
        font.pixelSize: 15
    }
}
\end{qml}

\begin{itemize}
  \item \cod{import QtQuick} trae los tipos básicos (\cod{Item}, \cod{Rectangle},
  \cod{Text}, \cod{Image}…). \cod{import app.musicplayer} trae los del proyecto, incluidos los
  objetos escritos en Rust (\cod{Player}, \cod{Library}…).
  \item \cod{id} es un nombre interno del archivo. No es una propiedad: no se puede cambiar.
  \item Los hijos se declaran dentro de las llaves del padre y se dibujan encima de él, en
  orden (el último, más arriba).
\end{itemize}

\section{Propiedades y \emph{bindings}}

La idea más importante de QML: cuando escribes una expresión como valor de una propiedad,
\textbf{no se evalúa una vez, se mantiene}. Si algo de lo que depende cambia, se recalcula.

\begin{qml}
Text {
    text: Player.playing ? "Pausa" : "Reproducir"
    color: Player.error !== "" ? Theme.accent : Theme.text
}
\end{qml}

Cuando Rust cambia \cod{Player.playing}, el texto cambia solo. Por eso el núcleo en Rust nunca
«dibuja»: solo cambia propiedades, y la interfaz las sigue.

\subsection{Propiedades propias}

\begin{qml}
Item {
    id: root
    property string titulo: ""                        // propiedad con valor inicial
    property int indice: 0
    property var elementos: []                        // var: cualquier valor de JavaScript
    readonly property bool vacio: elementos.length === 0   // calculada, de solo lectura
}
\end{qml}

Tipos que verás: \cod{int}, \cod{real}, \cod{bool}, \cod{string}, \cod{color}, \cod{var}
(objetos y listas de JavaScript) e \cod{Item} (una referencia a otro objeto visual).

\subsection{Romper un binding}

Si en código JavaScript \emph{asignas} un valor a una propiedad que tenía un binding
(\cod{root.indice = 3}), el binding desaparece y queda el valor fijo. Es intencionado y el
proyecto lo hace a menudo (por ejemplo, \cod{nowOpen} empieza ligado a un argumento de
depuración y luego se asigna al abrir o cerrar la pantalla «Sonando»).

\section{Colocar cosas}

\begin{qml}
Item {
    width: 400; height: 200

    Rectangle { x: 20; y: 10; width: 100; height: 100 }   // posición absoluta (relativa al padre)

    Text {
        anchors.right: parent.right                        // anclas: relativas a otros objetos
        anchors.verticalCenter: parent.verticalCenter
        anchors.rightMargin: 16
    }

    Row {                                                  // pone sus hijos en fila
        spacing: 8
        Rectangle { width: 30; height: 30 }
        Rectangle { width: 30; height: 30 }
    }
}
\end{qml}

El proyecto mezcla las tres formas: \cod{x}/\cod{y} calculados para las vistas animadas (la
estantería, la cubeta), \cod{anchors} para pegar cosas a bordes, y \cod{Row}/\cod{Column} para
grupos de botones y textos.

\section{Señales y manejadores}

Una \emph{señal} es un evento que un objeto emite; quien lo usa la escucha con un
\emph{manejador} \cod{on} + nombre con mayúscula:

\begin{qml}
// En IconButton.qml
Item {
    id: root
    signal clicked()
    TapHandler { onTapped: root.clicked() }
}

// Al usarlo
IconButton {
    icon: Theme.iconNext
    onClicked: Player.next()
}
\end{qml}

Toda propiedad tiene además una señal automática \cod{<propiedad>Changed}:

\begin{qml}
onIndexChanged: pos = index           // cada vez que cambia `index`
\end{qml}

Las señales pueden llevar datos; en el manejador se recogen con una función flecha:

\begin{qml}
SectionView { onActivated: item => win.choose(item) }
\end{qml}

\section{Funciones}

Se escriben en JavaScript dentro del objeto:

\begin{qml}
function step(delta) {
    if (items.length)
        index = Math.max(0, Math.min(items.length - 1, index + delta))
}
\end{qml}

Puedes llamar a funciones de otros objetos por su \cod{id} (\cod{sectionView.jump(3)}) y a las
funciones que expone Rust (\cod{Player.playIndex(2)}).

\section{Listas: \texttt{Repeater} y \texttt{ListView}}

\begin{qml}
ListView {
    width: 300; height: 400
    model: Catalog.queue                  // un array de JavaScript
    delegate: Text {
        required property var modelData  // el elemento de esta fila
        required property int index      // su posición
        text: (index + 1) + ". " + modelData.title
    }
}
\end{qml}

\begin{itemize}
  \item \cod{model} es la lista de datos; \cod{delegate}, la plantilla de cada fila.
  \item \cod{ListView} solo crea las filas visibles y se puede desplazar: para listas largas.
  \item \cod{Repeater} crea un objeto por elemento, sin desplazamiento: para listas cortas o
  cuando quieres colocarlos tú (los botones de pestaña, las cajas de la estantería).
\end{itemize}

\section{Interacción: los \emph{handlers}}

Qt Quick tiene dos generaciones de objetos para el ratón. El proyecto usa la moderna, los
\emph{pointer handlers}, salvo en casos concretos:

\begin{qml}
Item {
    HoverHandler { id: raton; cursorShape: Qt.PointingHandCursor }   // ¿está el ratón encima?
    TapHandler { onTapped: hacerAlgo(); onDoubleTapped: abrir() }   // clic, doble clic
    TapHandler { acceptedButtons: Qt.RightButton; onTapped: menu() }  // clic derecho
    WheelHandler { onWheel: event => desplazar(event.angleDelta.y) }  // rueda
}
\end{qml}

\begin{trampa}[Los clics que atraviesan]
Un \cod{TapHandler} normal no «se queda» el clic: lo que hay debajo también lo recibe. En los
botones y en las capas que tapan la pantalla hay que pedir que lo agarre en exclusiva con
\cod{gesturePolicy: TapHandler.ReleaseWithinBounds}. Lo verás en el
capítulo~\ref{cap:trampas}; ya nos pasó tres veces.
\end{trampa}

\section{Animaciones}

La forma más sencilla: \cod{Behavior}, que anima cualquier cambio de una propiedad.

\begin{qml}
Rectangle {
    x: abierto ? 0 : -300
    opacity: visible ? 1 : 0
    Behavior on x { NumberAnimation { duration: 320; easing.type: Easing.OutCubic } }
    Behavior on opacity { NumberAnimation { duration: 160 } }
}
\end{qml}

Cambias \cod{abierto} y el rectángulo se desliza. Las duraciones comunes del proyecto están en
\cod{Theme.animFast} (160 ms) y \cod{Theme.animMove} (320 ms). Para animaciones continuas
(el disco que gira) se usa \cod{RotationAnimation} con \cod{loops: Animation.Infinite}.

\section{Componentes reutilizables y \texttt{required}}

Cualquier archivo QML es un componente. Sus \cod{property} son su «interfaz pública»:

\begin{qml}
// PillButton.qml
Item {
    id: root
    property string text
    property bool primary: false
    signal clicked()
    // ...
}

// Uso
PillButton { text: "Aleatorio"; onClicked: root.play(0, true) }
\end{qml}

\section{Singletons}

Un \emph{singleton} es un objeto del que solo hay uno y que se usa por su nombre desde
cualquier archivo, sin crearlo. El proyecto tiene siete:

\begin{center}
\begin{tabular}{@{}lll@{}}
\toprule
Nombre & Escrito en & Archivo \\
\midrule
\cod{Player}, \cod{Library}, \cod{Account}, \cod{AppSettings}, \cod{Plugins} & Rust & \archivo{src/*.rs} \\
\cod{Theme}, \cod{Catalog} & QML & \archivo{qml/Theme.qml}, \archivo{qml/Catalog.qml} \\
\bottomrule
\end{tabular}
\end{center}

Un singleton QML empieza con \cod{pragma Singleton} y se registra en \archivo{build.rs} en
\cod{QML_SINGLETONS}.

\section{\texttt{Loader}: crear y destruir a demanda}

\begin{qml}
Loader {
    sourceComponent: vista === "lomos" ? lomos : (vista === "cubeta" ? cubeta : muro)
}
Component { id: lomos; LomosView { } }
\end{qml}

Un \cod{Loader} crea su contenido solo cuando hace falta y lo destruye al cambiar. El
proyecto lo usa para la vista de cada sección (\archivo{qml/SectionView.qml}) y para recrear
los iconos (lo verás en el capítulo~\ref{cap:trampas}).

\section{Efectos gráficos}

\begin{description}[leftmargin=1.2em]
  \item[\cod{MultiEffect}] (de \cod{QtQuick.Effects}): desenfoque, saturación, brillo,
  sombras y máscaras sobre cualquier elemento. Lo usan los paneles de vidrio, la letra
  desenfocada y los reflejos de la estantería.
  \item[\cod{layer.enabled}] Dibuja un elemento en una textura aparte (necesario para
  aplicarle efectos). Cuesta memoria de GPU: se activa solo donde hace falta.
  \item[\cod{ShaderEffectSource}] Copia en una textura lo que hay en otra parte de la escena:
  así los paneles de vidrio «ven» lo que tienen detrás para desenfocarlo.
  \item[\cod{Shape}] Dibujo vectorial (los iconos, a partir de trazos SVG).
  \item[\cod{Canvas}] Dibujo con una API como la de \texttt{<canvas>} de HTML. Solo se usa para
  leer los píxeles de una carátula pequeña y calcular su color dominante.
\end{description}

\section{Depurar QML}

\begin{itemize}
  \item \cod{console.log("valor", algo)} escribe en la terminal desde la que lanzaste el
  programa.
  \item Los errores de QML también salen ahí, con archivo y línea:
\begin{terminal}
qrc:/qt/qml/app/musicplayer/qml/LomosView.qml:312:25: Cannot override FINAL property
\end{terminal}
  \item Un error de sintaxis o de tipos en un componente hace que \emph{ese componente} no se
  cargue (y a menudo los que lo usan): si «desaparece» algo de la pantalla, mira la terminal.
\end{itemize}

\section{Para aprender más}

La documentación oficial de Qt es muy buena: \url{https://doc.qt.io/qt-6/qmlapplications.html}.
Lee en especial «QML Syntax Basics», «Property Binding», «Signal and Handler Event System» y
«Positioning with Anchors».

\begin{ejercicio}
Abre \archivo{qml/PillButton.qml} y léelo entero. Identifica: sus propiedades públicas, su
señal, el binding que decide el color de fondo y el manejador que emite la señal. Después haz
lo mismo con \archivo{qml/Slider.qml}, que es algo más complejo (arrastrar).
\end{ejercicio}
